%%%%%%%%%%%%%%%%%%%%%%%%%%%%%%%%%%%%%%%%%%%%%%%%%%%%%%%%%%%%%%%%%%%%%%%%%%%%%%%
%  a1954 12Be analysis -- step-by-step user guide
%  Build:  pdflatex a1954_Be12_guide.tex   (twice, for the ToC)
%%%%%%%%%%%%%%%%%%%%%%%%%%%%%%%%%%%%%%%%%%%%%%%%%%%%%%%%%%%%%%%%%%%%%%%%%%%%%%%
\documentclass[11pt,a4paper]{article}

\usepackage[T1]{fontenc}
\usepackage[utf8]{inputenc}
\usepackage[margin=2.3cm]{geometry}
\usepackage{amsmath}
\usepackage{booktabs}
\usepackage{longtable}
\usepackage{array}
\usepackage{xcolor}
\usepackage{listings}
\usepackage{enumitem}
\usepackage{fancyhdr}
\usepackage{caption}
\usepackage{microtype}
\usepackage[colorlinks=true,linkcolor=blue!55!black,urlcolor=blue!55!black,
            citecolor=blue!55!black]{hyperref}

% ---- literal code (never math mode; underscores stay literal) ------------------
\newcommand{\code}[1]{\texttt{\detokenize{#1}}}
\newcommand{\file}[1]{\texttt{\detokenize{#1}}}

% ---- shell / macro listings ---------------------------------------------------
\definecolor{shellbg}{RGB}{245,245,242}
\definecolor{shellrule}{RGB}{200,200,192}
\definecolor{cmt}{RGB}{100,120,100}
\lstdefinestyle{sh}{
  basicstyle=\ttfamily\footnotesize,
  backgroundcolor=\color{shellbg},
  frame=single, rulecolor=\color{shellrule},
  framesep=5pt, xleftmargin=6pt, xrightmargin=2pt,
  breaklines=true, breakatwhitespace=false,
  postbreak=\mbox{\textcolor{gray}{$\hookrightarrow$}\space},
  columns=fullflexible, keepspaces=true,
  commentstyle=\color{cmt}\itshape,
  morecomment=[l]{\#}, showstringspaces=false,
  literate={\_}{{\_}}1 {~}{{\textasciitilde}}1
}
\lstset{style=sh}

\pagestyle{fancy}
\fancyhf{}
\lhead{\small a1954 $^{12}$Be --- analysis guide}
\rhead{\small\thepage}
\renewcommand{\headrulewidth}{0.4pt}

\newcommand{\warn}[1]{\par\noindent\fcolorbox{red!45!black}{red!4}{%
  \parbox{\dimexpr\linewidth-2\fboxsep-2\fboxrule}{\small #1}}\par\medskip}
\newcommand{\note}[1]{\par\noindent\fcolorbox{blue!40!black}{blue!4}{%
  \parbox{\dimexpr\linewidth-2\fboxsep-2\fboxrule}{\small #1}}\par\medskip}

\title{\vspace{-1.4cm}\textbf{Running the a1954 $^{12}$Be analysis}\\[2mm]
\large Step-by-step guide to the reconstruction, pattern-recognition and\\
track-fitting pipeline, with the tuning knobs that matter}
\author{ATTPCROOT \texttt{OpenKF-Claude} branch}
\date{\today}

\begin{document}
\maketitle
\thispagestyle{fancy}

\begin{abstract}
\noindent
This guide walks through the complete analysis of the a1954 $^{12}$Be proton-target
data set --- $^{12}$Be$(p,p')^{12}$Be$^{*}$ and $^{12}$Be$(p,d)^{11}$Be in inverse
kinematics --- from the raw HDF5 files to the excitation-energy spectrum. Every stage is
given as a command you can copy, followed by the options that stage exposes. Sections
\ref{sec:pra-tuning} and \ref{sec:fit-tuning} collect the pattern-recognition (PRA) and
fitting parameters in reference tables together with recipes for improving them, and
Section~\ref{sec:pitfalls} lists the traps this data set has already sprung.
\end{abstract}

\tableofcontents
\bigskip

%===============================================================================
\section{What the analysis does}
%===============================================================================

A $^{12}$Be beam at $\approx 155$~MeV (12.9~MeV/u, \emph{not yet independently confirmed};
see Section~\ref{sec:pitfalls}) is stopped in the AT-TPC filled with H$_2$ at 600~torr,
which acts simultaneously as target and detection medium. Two channels are analysed:

\begin{center}
\begin{tabular}{@{}lllll@{}}
\toprule
Channel & Ejectile detected & Typical $\theta_{\rm lab}$ & Typical KE & Residual \\
\midrule
$^{12}$Be$(p,p')^{12}$Be$^{*}$ & recoil proton   & $60$--$90^\circ$ (backward) & 1--10 MeV  & $^{12}$Be \\
$^{12}$Be$(p,d)^{11}$Be        & deuteron        & $10$--$50^\circ$ (forward)  & 10--30 MeV & $^{11}$Be \\
\bottomrule
\end{tabular}
\end{center}

The pipeline is staged so that the expensive part (unpacking + PSA + PRA, hours per run) is
done \emph{once} and the cheap part (fitting, minutes per run) can be iterated freely:

\begin{center}\small
\file{run_XXXX.h5} $\rightarrow$ \textbf{reco} (unpack + PSA + clean + PRA) $\rightarrow$
\file{run_XXXX_reco.root} $\rightarrow$ \textbf{slim} $\rightarrow$ \textbf{gate}
(IC + PID) $\rightarrow$ \textbf{fit} (UKF and/or GENFIT) $\rightarrow$
\textbf{kinematics} ($E_x$ spectrum, explorers)
\end{center}

\noindent All macros live in
\file{macro/Unpack_HDF5/a1954_Be12/UKF/} of the \code{OpenKF-Claude} branch, referred to
below as \textbf{\code{$HERE}}.

%===============================================================================
\section{Setting up}
%===============================================================================

\subsection{Environment}

\begin{lstlisting}
export REPO=/home/yassid/fair_install/ATTPCROOTv2-OpenKF
export HERE=$REPO/macro/Unpack_HDF5/a1954_Be12/UKF

source /home/yassid/fair_install/FairSoft/install/bin/thisroot.sh
source $REPO/build/config.sh
# MANDATORY -- config.sh does not set this, and cling crashes without it
export ROOT_INCLUDE_PATH=$REPO/build/include:/home/yassid/fair_install/FairRootInstall/include
\end{lstlisting}

\warn{\textbf{\code{ROOT_INCLUDE_PATH} must point at \file{build/include}}, not at the
source directories. Pointing it at the sources makes cling crash on a hidden-friend
\code{swap} declaration. Every batch script in \code{$HERE} already exports it correctly;
you only need this line for interactive work.}

\subsection{Two builds, on purpose}

\begin{center}
\begin{tabular}{@{}lll@{}}
\toprule
Build & GENFIT2 & Use for \\
\midrule
\file{build/}        & not found & reco, PSA, PRA, UKF fitting \\
\file{build_genfit/} & linked against \file{~/fair_install/GenFitInst} & \code{fitGenfit_Be12.C} \\
\bottomrule
\end{tabular}
\end{center}

\noindent They compile the same sources; they are kept apart so a GENFIT rebuild never
disturbs a running reco batch. Source \file{build_genfit/config.sh} instead of
\file{build/config.sh} when you run GENFIT.

\subsection{Geometry (once per clone)}

The geometry \code{.root} files are \emph{not} tracked in git. Build the 600~torr
hydrogen geometry that GENFIT needs before the first GENFIT fit:

\begin{lstlisting}
cd $REPO && root -b -q geometry/ATTPC_H600torr.C
# writes geometry/ATTPC_H600torr.root and ATTPC_H600torr_geomanager.root
\end{lstlisting}

\warn{Do \emph{not} fall back on \file{ATTPC_H1bar_geomanager.root}: a1954 ran at 600~torr,
so the 1-bar geometry is 27\,\% too dense and every material-effects correction GENFIT
applies will be miscalibrated.}

\subsection{Data}

\begin{center}
\begin{tabular}{@{}lp{10.2cm}@{}}
\toprule
Raw HDF5      & \file{/mnt/f/a1954_remerged/run_XXXX.h5} (drive F, \code{WD_Cris}) \\
Mount         & \code{sudo mount -t drvfs F: /mnt/f} --- F: is \emph{not} auto-mounted after a reboot \\
Physics runs  & 0142--0215 (+0216--0224); 0160 and 0177--0181, 0183 are absent \\
Runs fitted   & 75 (\file{~/a1954_Be12_fit/}); \textbf{73 clean} after dropping the two desynced runs \\
Excluded      & \textbf{run\_0142} and \textbf{run\_0212} (timestamp desync, Section~\ref{sec:tsync});
                \textbf{run\_0161} in the $(p,d)$ pass (FRIB file empty, 587~B) \\
Parameters    & \file{parameters/ATTPC.a1954_Be12.par} \\
\bottomrule
\end{tabular}
\end{center}

\noindent Key entries of the parameter file:

\begin{center}
\begin{tabular}{@{}llp{7.4cm}@{}}
\toprule
Parameter & Value & Comment \\
\midrule
\code{BField}        & 2.85 T   & sign is set in the macros, \emph{not} here \\
\code{GasPressure}   & 600 torr & H$_2$ \\
\code{DriftVelocity} & 1.00 cm/$\mu$s & \emph{the} parameter to scan (Section~\ref{sec:driftv}) \\
\code{SamplingRate}  & 3.0      & MHz \\
\code{TBEntrance}    & 300      & beam entrance in time buckets \\
\code{NumTbs}        & 512      & \\
\code{ZPadPlane}     & 1000 mm  & \\
\code{PeakingTime}   & 720 ns   & GET electronics $\tau$ \\
\bottomrule
\end{tabular}
\end{center}

\note{The $^{14}$C sibling analysis (\file{macro/Unpack_HDF5/a1954/UKF/}) uses a
\emph{different} par file with \code{SamplingRate 6.0}, \code{DriftVelocity 1.30} and
\code{TBEntrance 480}. Do not mix them up.}

%===============================================================================
\section{Stage 1 --- reconstruction (unpack $+$ PSA $+$ cleaning $+$ PRA)}
\label{sec:reco}
%===============================================================================

This is the expensive stage: it reads the 5--9~GB HDF5 file, turns pad traces into hits,
removes noise hits and clusters the surviving hits into track candidates.

\subsection{One run, interactively}

\begin{lstlisting}
root -b -q -l "$HERE/pipeline/unpackReco_Be12.C(\"run_0147\")"
# first 500 events only, for a quick look:
root -b -q -l "$HERE/pipeline/unpackReco_Be12.C(\"run_0147\",500)"
\end{lstlisting}

\subsection{Production (the validated configuration)}

\begin{lstlisting}
cd $HERE
./reco_hdb_F.sh "run_0142 run_0143 run_0144" 2      # 2-parallel; drvfs I/O is the cap
\end{lstlisting}

\noindent \code{reco_hdb_F.sh} reads from \file{/mnt/f/a1954_remerged/} and writes\\
\file{/mnt/f/a1954_Be12_reco_hdb/run_XXXX_reco.root}. It skips runs that are already
done and logs to \file{$OUT/logs/}. The configuration it pins is:

\begin{center}
\textbf{AtPSAMultiFit (thr 20) $+$ AtDirDeDxCleaner $+$ AtTrackFinderHDBSCAN (mover join)}
\end{center}

\subsection{Full argument list of \texttt{unpackReco\_Be12.C}}

\begin{longtable}{@{}lll p{6.9cm}@{}}
\toprule
\# & Argument & Default & Meaning \\
\midrule
\endhead
1  & \code{fileName}      & \code{"run_0142"} & run name, no extension \\
2  & \code{nEvents}       & 0                 & 0 = all events \\
3  & \code{persistRaw}    & \code{true}       & write the raw pads too; production uses \code{false} (files are 4$\times$ smaller) \\
4  & \code{outDir}        & \file{~/a1954_Be12_reco/} & output directory \\
5  & \code{filepath}      & (NSCL drive)      & input directory; use \file{/mnt/f/a1954_remerged/} \\
6  & \code{doSC}          & \code{false}      & space-charge correction; not needed (does not affect $z$) \\
7  & \code{applyTimeCorr} & \code{false}      & pad-time offsets --- the CSV that exists is a1975-specific; leave OFF until an a1954 table is measured \\
8  & \code{psaType}       & \code{"multifit"} & \code{"multifit"} or \code{"max"} \\
9  & \code{primSigma}     & 0                 & $>0$ enables prominence-on-primary at that many $\sigma$ \\
10 & \code{thr}           & 20                & PSA amplitude threshold [ADC] \\
11 & \code{praType}       & \code{"tc"}       & \code{"tc"} (triplclust) or \code{"hdbscan"} --- \textbf{production uses hdbscan} \\
12 & \code{hdMcs}         & 20                & HDBSCAN min cluster size \\
13 & \code{hdMs}          & 8                 & HDBSCAN min samples \\
14 & \code{fitChi2}       & 0                 & $>0$ = PSA shape $\chi^2$ gate (also turns on float-$\tau$); $<0$ = diagnostic only \\
15 & \code{relErr}        & 0.1               & fractional model error in the PSA $\chi^2$ \\
16 & \code{parFile}       & \code{ATTPC.a1954_Be12.par} & \\
17 & \code{doClean}       & \code{kTRUE}      & run \code{AtDirDeDxCleaner} before PRA \\
\bottomrule
\end{longtable}

\subsection{What each sub-stage produces}

\begin{center}
\begin{tabular}{@{}llp{7.3cm}@{}}
\toprule
Task & Branch out & Content \\
\midrule
\code{AtUnpackTask}  & \code{AtRawEvent} & raw pad traces (only if \code{persistRaw}) \\
\code{AtPSAtask}     & \code{AtEventH}   & \code{AtHit} point cloud \\
\code{AtDataCleaningTask} & \code{AtEventClean} & de-noised point cloud (transient) \\
\code{AtPRAtask}     & \code{AtPatternEvent} & \code{AtTrack} candidates + their hits \\
\bottomrule
\end{tabular}
\end{center}

%===============================================================================
\section{Stage 2 --- FRIB/ion-chamber unpacking and the slim cache}
%===============================================================================

The ion chamber lives in the \code{frib/evt} HDF5 group, not in the pad stream, and is
unpacked separately by \code{AtFRIBHDFUnpacker}. It is fast (no PSA/PRA).

\begin{lstlisting}
cd $HERE
./frib_F_batch.sh "run_0142 run_0143 ..." 4   # -> ~/a1954_Be12_reco_hdb_slim/run_XXXX_FRIB.root
./slim_cache_batch.sh "run_0142 run_0143 ..." 4  # -> ~/a1954_Be12_reco_hdb_slim/run_XXXX_slim.root
\end{lstlisting}

\noindent The \emph{slim cache} clones only the \code{AtPatternEvent} and
\code{EventHeader} branches of the big reco file onto the local disk. Every downstream
PID/gate/fit pass then reads locally instead of streaming tens of GB back over drvfs
($\approx 250\times$ faster).

\warn{An empty FRIB file (a few hundred bytes, e.g. run\_0161) means no IC gate is
possible for that run. The batch scripts skip files below 10~kB.}

%===============================================================================
\section{Stage 3 --- beam identification and particle ID}
\label{sec:pid}
%===============================================================================

\subsection{The ion-chamber gate}

The a1954 beam is a fragment cocktail. The IC pulse is \emph{bipolar}: a positive
overshoot near TB~1134 and a negative dip near TB~1155.

\begin{center}
\begin{tabular}{@{}ll@{}}
\toprule
IC observable & \textbf{maximum} ADC in TB $[1050,1250]$ (the positive overshoot, \emph{not} $-\min$) \\
$^{12}$Be     & the isolated peak at $\approx 650$ $\Rightarrow$ gate \textbf{[500, 900]} \\
Contaminant   & the dominant peak at $\approx 1900$ --- reject \\
Pile-up       & require \textbf{exactly one} pulse above 200 ADC in TB $[800,1500]$ \\
\bottomrule
\end{tabular}
\end{center}

\begin{lstlisting}
root -b -q -l "$HERE/pid/ic_Be12.C(\"run_0147,run_0148\",\"~/a1954_Be12_reco_hdb_slim/\")"
root -b -q -l "$HERE/pid/ic_multiplicity_Be12.C(\"run_0147\")"   # pile-up check
\end{lstlisting}

\subsection{Timestamp synchronisation --- run this before trusting any gate}
\label{sec:tsync}

The IC is matched to the pad data \textbf{by entry index}, which is only valid when the two
streams have not dropped events relative to each other. The check compares timestamp
\emph{deltas} (so it is offset-independent):

\begin{lstlisting}
root -b -q -l "$HERE/pid/tsync_check.C(\"run_0143\")"
\end{lstlisting}

\begin{center}
\begin{tabular}{@{}lll@{}}
\toprule
Verdict & Signature & Action \\
\midrule
aligned  & event-count difference $+1$ (trailing event), $\approx 100\,\%$ deltas agree & use \\
desynced & negative count difference, $0\,\%$ agree & \textbf{exclude the run} \\
\bottomrule
\end{tabular}
\end{center}

\noindent For a1954 $^{12}$Be: 73 of 75 runs are aligned; \textbf{run\_0142 (diff $-9$)} and
\textbf{run\_0212 (diff $-5$)} are desynced and must be dropped.

\subsection{Drawing the PID gate}

The PID plane is $\sqrt{dE/dx}$ versus $B\rho$ from \code{AtTools::AtSpyralPID}, computed on
the PRA track candidates --- \emph{no fit is needed}, so this is the natural place to
inspect the data before committing to a fitting configuration. It is also charge-sign
agnostic, hence unaffected by the $B_z$ handedness.

\begin{lstlisting}
# 1. build the PID plane and persist it as a CSV
root -b -q -l "$HERE/pid/pid_Be12.C(\"run_0147,run_0148\",\"~/a1954_Be12_reco_hdb_slim/\")"
root -b -q -l "$HERE/pid/dump_pid_Be12.C(\"run_0147\",\"<recoDir>\",\"$HERE/pid/pid_Be12_data.csv\")"

# 2. draw the polygon by hand (INTERACTIVE -- never with -b)
root -l "$HERE/pid/draw_gate_Be12.C(\"proton_12Be\")"
root -l "$HERE/pid/draw_gate_Be12.C(\"deuteron_12Be\",625,750,0,1,2)"

# 3. verify it headlessly
root -b -q -l "$HERE/pid/apply_gate_Be12.C(\"proton_12Be.json\")"
root -b -q -l "$HERE/pid/show_gates_Be12.C"
\end{lstlisting}

\noindent Left-click to place each vertex, double-click to close; the polygon is written to
\file{<name>.json} in \file{$HERE/pid/}. The gates in use are
\file{pid/proton_12Be.json} (wraps the full proton band) and \file{pid/deuteron_12Be.json}.

\note{An \code{arclen > 200} cut visibly sharpens the band and is worth trying when
redrawing a gate. HDBSCAN reco gives a tighter band than triplclust.}

%===============================================================================
\section{Stage 4 --- building the gated fit input}
%===============================================================================

\code{gate_events_Be12.C} writes a small ROOT file containing \emph{only} the events that
pass the IC gate and, inside them, \emph{only} the tracks that fall in the PID polygon and
the requested angular window. Because it also copies the FairRoot metadata
(\code{BranchList}, \code{FileHeader}, \code{cbmout}), the fitters read it exactly like a
normal reco file.

\begin{lstlisting}
root -b -q -l "$HERE/pipeline/gate_events_Be12.C(\"run_0147\",\
  \"/home/yassid/a1954_Be12_reco_hdb_slim/\",\"/home/yassid/a1954_Be12_fit/in/\",\
  \"/mnt/f/a1954_Be12_reco_hdb/run_0147_reco.root\")"
\end{lstlisting}

\begin{longtable}{@{}lll p{6.6cm}@{}}
\toprule
\# & Argument & Default & Meaning \\
\midrule
\endhead
1  & \code{run}         & ---   & run name \\
2  & \code{inDir}       & slim dir & holds \file{_slim.root} and \file{_FRIB.root} \\
3  & \code{outDir}      & \file{~/a1954_Be12_fit/in/} & gated output \\
4  & \code{refReco}     & \code{""} & big reco file, only to copy FairRoot metadata from \\
5,6 & \code{icLo, icHi} & 500, 900 & IC amplitude window \\
7,8 & \code{icTbLo, icTbHi} & 1050, 1250 & IC integration window [TB] \\
9  & \code{bField}      & 2.85  & used by \code{AtSpyralPID} (magnitude only) \\
10 & \code{protonGate}  & \file{pid/proton_12Be.json} & PID polygon; swap for \file{deuteron_12Be.json} \\
11 & \code{thMin}       & 90.0  & minimum $\theta_{\rm lab}$ [deg]. \textbf{90 for $(p,p')$; 0 for $(p,d)$} \\
12 & \code{peakThr}     & 200   & IC pulse-counting threshold \\
13,14 & \code{pkTbLo, pkTbHi} & 800, 1500 & IC pulse-counting window \\
\bottomrule
\end{longtable}

\warn{The $\theta$ convention differs between stages. The \code{thMin} cut here is on the
\emph{Spyral polar} angle \code{r.polar}, whereas the fitted \code{k.theta} used downstream
satisfies $\theta_{\rm fit} \approx 180^\circ - \theta_{\rm polar}$. Using
\code{thMin = 90} on a $(p,d)$ pass kills every deuteron.}

%===============================================================================
\section{Stage 5 --- fitting}
\label{sec:fit}
%===============================================================================

Both fitters read the \emph{same} gated input, so tracks can be compared one-to-one.

\subsection{The $(p,p')$ production pass}

\begin{lstlisting}
cd $HERE
./fitpipe_Be12.sh "run_0143 run_0144 run_0145" 4     # gate -> UKF -> GENFIT, 4 cores
# outputs: ~/a1954_Be12_fit/run_XXXX_{ukf,genfit}.root  (+ rsync to /mnt/f/a1954_Be12_fit/)
\end{lstlisting}

\subsection{The $(p,d)$ production pass}

\begin{lstlisting}
./fitpipe_pd_Be12.sh "run_0143 run_0144" 4 0   # deuteron gate, thMin = 0 (FORWARD tracks)
./genfit_pd_Be12.sh  "run_0143 run_0144" 4     # GENFIT over the same gated input
# outputs: ~/a1954_Be12_fit_pd/  and  ~/a1954_Be12_genfix_pd/
\end{lstlisting}

\subsection{UKF --- \texttt{fitUKF\_Be12.C}}

\begin{lstlisting}
root -b -q -l "$HERE/pipeline/fitUKF_Be12.C(\"run_0147\",-1,\"proton\",-1,2.85,6.5e-5,\
  \"\",\"/home/yassid/a1954_Be12_fit/in/\",0.5,0.1,1,10,\"/home/yassid/a1954_Be12_fit/\")"
\end{lstlisting}

\begin{longtable}{@{}ll l p{6.3cm}@{}}
\toprule
\# & Argument & Default & Meaning \\
\midrule
\endhead
1  & \code{fileName}      & \code{"run_0142"} & \\
2  & \code{nEvents}       & $-1$ & $-1$ = all \\
3  & \code{particles}     & \code{"proton"} & comma list; one hypothesis per species (\code{proton}, \code{deuteron}, \code{triton}, \code{He3}, \code{alpha}) \\
4  & \code{bFieldSign}    & $-1$ & \textbf{$-1$ = experimental handedness}, $+1$ = simulation \\
5  & \code{bFieldMag}     & 2.85 & T \\
6  & \code{gasDensity}    & 6.5e-5 & g/cm$^3$, feeds CATIMA \\
7  & \code{outSuffix}     & \code{""} & appended to the output name (use for scans) \\
8  & \code{ioDir}         & reco dir & input directory \\
9  & \code{measSigma}     & 0.5 & \textbf{per-cluster position $\sigma$ [mm]} \\
10 & \code{momSigmaFrac}  & 0.1 & fractional momentum-seed uncertainty \\
11 & \code{nIter}         & 1   & filter/smoother iterations \\
12 & \code{minClusters}   & 10  & minimum clusters for a fit attempt \\
13 & \code{outDir}        & \code{""} & defaults to \code{ioDir} \\
14 & \code{refTrack}      & \code{false} & reference-track mode \\
15 & \code{refInflation}  & 4.0 & covariance inflation in reference-track mode \\
\bottomrule
\end{longtable}

\warn{\textbf{Handedness.} Experimental data has the opposite helix handedness to
simulation because the digitisation reverses the drift $z$. Use $B_z = -2.85$~T for data.
This was validated on run\_0142: median $\chi^2/\mathrm{ndf} = 0.03$, median KE 5.96~MeV,
clean elastic locus. If fits diverge or the KE is unphysical, flip the sign first.}

\subsection{GENFIT --- \texttt{fitGenfit\_Be12.C}}

\begin{lstlisting}
source $REPO/build_genfit/config.sh          # <-- the GENFIT build
root -b -q -l "$HERE/pipeline/fitGenfit_Be12.C(\"run_0147\",-1,\"<inDir>\",\"\",\"<outDir>\",\
  -2.85,2,5,\"\",4.0,10.0,170.0,kTRUE,kFALSE,\"proton\",\"ATTPC_H600torr\")"
\end{lstlisting}

\begin{longtable}{@{}ll l p{6.1cm}@{}}
\toprule
\# & Argument & Default & Meaning \\
\midrule
\endhead
1  & \code{fileName}       & \code{"run_0142"} & \\
2  & \code{nEvents}        & $-1$ & \\
3  & \code{ioDir}          & reco dir & \\
4  & \code{outSuffix}      & \code{""} & \\
5  & \code{outDir}         & \code{""} & \\
6  & \code{bField}         & $-2.85$ & signed, as for the UKF \\
7,8 & \code{minIter, maxIter} & 2, 5 & Kalman refit iterations \\
9  & \code{pidGate}        & \code{""} & optional in-fitter Spyral PID gate (JSON) \\
10 & \code{measSigma}      & 4.0 & \textbf{per-cluster position $\sigma$ [mm]} \\
11,12 & \code{thetaMinDeg, thetaMaxDeg} & 10, 170 & \textbf{lower it to 5 for $(p,d)$} \\
13 & \code{matEffects}     & \code{kFALSE} & \textbf{set \code{kTRUE}} --- see below \\
14 & \code{backwardSeedFix}& \code{kFALSE} & alternative backward-track seeding \\
15 & \code{particle}       & \code{"proton"} & \code{proton}, \code{deuteron}, \code{triton} \\
16 & \code{geoName}        & \code{ATTPC_H600torr} & must match the real gas density \\
\bottomrule
\end{longtable}

\note{\textbf{Always run GENFIT with \code{matEffects = kTRUE} and the 600-torr geometry.}
With \code{matEffects = kFALSE} GENFIT fits a constant-momentum helix, so the reported KE is
a track \emph{average}: it comes out $\approx 30\,\%$ low below 6~MeV, 9.5\,\% low at
6--10~MeV, and only converges to the UKF above $\approx 10$--15~MeV. That single flag was
the whole $+0.45$~MeV $E_x$ offset.}

\subsection{Which fitter to believe}

After the \code{AtGenfitter} energy-loss fix (commit \code{5f87e625}) the two fitters
agree. The bug: the code called \code{useEnergyLossParam()} unconditionally, which threw on
every step and silently truncated reference tracks to $\mathrm{ndf} \le 0$.

\begin{center}
\begin{tabular}{@{}lrrr@{}}
\toprule
$(p,p')$, 73 clean runs, $\chi^2/\mathrm{ndf} < 5$ & tracks & g.s.\ $\mu$ [MeV] & FWHM [MeV] \\
\midrule
UKF                          & 26\,604 & $+0.020$ & 0.555 \\
GENFIT, pre-fix (\code{matEffects=kFALSE}) & 12\,533 & $+0.451$ & --- \\
GENFIT, fixed                & 23\,316 & $+0.028$ & 0.521 \\
\bottomrule
\end{tabular}
\end{center}

\noindent GENFIT is now marginally the better resolution and its old KE-vs-$\theta$ double
band is gone. A known residual remains: a $+16\,\%$ KE deficit \emph{below 2~MeV}, because
under $\beta\gamma = 0.05$ genfit leaves Bethe--Bloch for \code{dEdxParam}, which returns
zero without a table. Load a proton $dE/dx$ table via \code{SetELossFile()} to close it.

For $(p,d)$: UKF 2\,661 vs GENFIT 2\,866 deuterons; both behave, because forward deuterons
are 10--30~MeV where GENFIT was never in trouble.

%===============================================================================
\section{Stage 6 --- physics output}
%===============================================================================

\subsection{Excitation-energy spectrum}

\code{ex_Be12.C} applies relativistic two-body kinematics to the fitted (KE, $\theta$) and
fills the $E_x$ spectrum, $E_x$ vs $\theta_{\rm cm}$, and KE vs $\theta_{\rm lab}$. It also
writes an \code{TNtuple} cache (\code{ke:theta:vertexz:thcm:ex:chi2ndf}) that every
interactive tool then works from.

\begin{lstlisting}
RUNS="run_0143,run_0144,run_0145"   # comma separated!

# (p,p')  -- proton ejectile, 12Be residual
root -b -q -l "$HERE/pp/ex_Be12.C(\"$RUNS\",\"/home/yassid/a1954_Be12_genfix/\",155,5,\
  \"_pp_genfit\",1.007825,12.026921,\"12Be(p,p')\",\"genfit\")"

# (p,d)   -- deuteron ejectile, 11Be residual
root -b -q -l "$HERE/pp/ex_Be12.C(\"$RUNS\",\"/home/yassid/a1954_Be12_fit_pd/\",155,5,\
  \"_pd_ukf\",2.014102,11.021658,\"12Be(p,d)11Be\",\"ukf\")"
\end{lstlisting}

\begin{center}
\begin{tabular}{@{}lll@{}}
\toprule
\# & Argument & Meaning \\
\midrule
1 & \code{runsCSV}   & comma-separated run list \\
2 & \code{inDir}     & directory holding \file{<run>_<fitter>.root} \\
3 & \code{Ebeam}     & beam energy [MeV] --- \textbf{155 for a1954 $^{12}$Be} \\
4 & \code{chi2Cut}   & $\chi^2/\mathrm{ndf}$ cut (5 nominal; \code{1e9} to disable) \\
5 & \code{outTag}    & cache suffix --- \textbf{must include the leading underscore} \\
6 & \code{mEjectAmu} & 1.007825 (p) or 2.014102 (d) \\
7 & \code{mResidAmu} & 12.026921 ($^{12}$Be) or 11.021658 ($^{11}$Be) \\
8 & \code{chanLabel} & plot title \\
9 & \code{fitter}    & \code{"ukf"} or \code{"genfit"} --- selects the input file suffix \\
\bottomrule
\end{tabular}
\end{center}

\noindent Caches land in \file{$HERE/pp/plots/proton_kin<outTag>.root}.

\subsection{Interactive exploration}

Both explorers recompute $E_x$ and $\theta_{\rm cm}$ live from the cache, so you can move
$E_{\rm beam}$, the cuts and the binnings without refitting anything. Both carry the UKF
and GENFIT caches with an in-page switch and an overlay.

\begin{lstlisting}
# browser version (no X11 needed; regenerates the HTML and launches Chrome)
$HERE/pp/open_explorer.sh pp        # 12Be(p,p')
$HERE/pp/open_explorer.sh pd        # 12Be(p,d)11Be

# ROOT/X11 GUI version -- plain ROOT, no config.sh needed, and never with -b
root -l $HERE/pp/explore_Be12.C
\end{lstlisting}

\warn{ROOT's 3D/Eve viewers crash under WSL, and \code{explorer.exe <file>} silently does
nothing on this machine --- \code{open_explorer.sh} therefore stages the HTML on the Windows
side and launches the browser binary directly. Prefer the browser explorer.}

\subsection{Fitter comparison}

\begin{lstlisting}
root -b -q -l "$HERE/pp/compare_fitters_Be12.C(\"run_0143,run_0144\",\
  \"/home/yassid/a1954_Be12_fit/\")"
# -> pp/plots/compare_fitters_Be12.png, compare_fitters_chi2_Be12.png
\end{lstlisting}

%===============================================================================
\section{Tuning the pattern recognition}
\label{sec:pra-tuning}
%===============================================================================

\warn{\textbf{PRA and fitting settings are physics-dependent.} A knob that improves one
reaction/geometry can hurt another. Always re-validate on \emph{this} data set --- one run,
a few hundred events, is enough to see the effect.}

\subsection{Choosing the track finder}

\begin{center}
\begin{tabular}{@{}lp{5.2cm}p{5.2cm}@{}}
\toprule
& \code{AtTrackFinderTC} (triplclust) & \code{AtTrackFinderHDBSCAN} \\
\midrule
Selected by & \code{praType="tc"} & \code{praType="hdbscan"} \\
Strength & fast, robust on isolated tracks & separates \emph{overlapping/crossing} tracks that triplclust merges \\
Weakness & merges crossing tracks & over-segments; needs the join step \\
On a1954 & tighter-band alternative & \textbf{production choice}; gives the tightest PID band \\
\bottomrule
\end{tabular}
\end{center}

\subsection{HDBSCAN parameters}

Set inside \code{unpackReco_Be12.C}; the first two are exposed as macro arguments
(\code{hdMcs}, \code{hdMs}), the rest are edited in the macro body.

\begin{longtable}{@{}l l l p{6.0cm}@{}}
\toprule
Setter & a1954 value & Default & Effect / how to tune \\
\midrule
\endhead
\code{SetMinClusterSize}        & 20   & adaptive & Smallest accepted cluster. \textbf{Lower} to recover short/low-energy tracks, \textbf{raise} to suppress noise blobs. Primary yield knob. \\
\code{SetMinSamples}            & 8    & 3  & $k$ for the core distance. Higher = more conservative, more points declared noise. \\
\code{SetClusterSelectionEpsilon} & 10.0 & 10.0 & [mm] Do not split clusters closer than this. \textbf{The big over-segmentation lever}: raise to stop one track being cut into pieces, lower to separate two nearby tracks. \\
\code{SetJoinMethod}            & \code{"mover"} & \code{"overlap"} & \code{mover} = vertex-seeded trajectory following; \code{overlap} = circle-area overlap; \code{continuity} = end-to-end; \code{none}. \\
\code{SetMinClusterSizeJoin}    & 15   & 15 & Minimum size a fragment must have to take part in a join. \\
\code{SetCircleOverlapRatio}    & 0.25 & 0.25 & Overlap join: merge if overlap area $>$ ratio $\times$ smaller area. \\
\code{SetMotionGapTol}          & 40   & 40 & [mm] \code{mover}: largest end-to-end gap that can be chained. Raise if tracks stay broken across a gap. \\
\code{SetMotionAngleTol}        & 35   & 35 & [deg] \code{mover}: tangent-continuity tolerance at the junction. \\
\code{SetOutlierScaleFactor}    & --- (0.05) & 0.05 & Per-cluster LOF cleaning; $\le 0$ disables. \\
\code{SetOutlierThreshold}      & --- (1.5)  & 1.5  & LOF above this $\rightarrow$ noise. \\
\code{SetJoinZFraction}         & --- (0.2)  & 0.2  & \code{continuity} join only. \\
\code{SetJoinRadiusFraction}    & --- (0.3)  & 0.3  & \code{continuity} join only. \\
\bottomrule
\end{longtable}

\subsection{Triplclust parameters}

Exposed through \code{SetScluster} ($s$, 0.3), \code{SetKtriplet} ($k$, 19),
\code{SetNtriplet} ($n$, 2), \code{SetMcluster} ($m$, 15), \code{SetRsmooth} ($r$, 2),
\code{SetAtriplet} ($a$, 0.03), \code{SetTcluster} ($t$, 4.0); a1954 additionally sets the
base-class \code{SetClusterRadius(15.0)} and \code{SetClusterDistance(7.5)}. Increasing $k$
and $m$ makes clustering more conservative (fewer, cleaner tracks); increasing $s$ merges
more aggressively.

\subsection{Base-class knobs shared by both finders (\code{AtPRA})}

\begin{center}
\begin{tabular}{@{}l l p{8.5cm}@{}}
\toprule
Setter & Default & Effect \\
\midrule
\code{SetMinHits} / \code{SetMaxHits} & 0 / 5000 & accepted hit-count window per track \\
\code{SetClusterRadius}   & 0 (15 for tc) & [mm] radius of the hit clusters handed to the fitter \\
\code{SetClusterDistance} & 0 (7.5 for tc)& [mm] spacing between those clusters \\
\code{SetTargetClusters}  & 25 & adaptive clustering target \\
\code{SetUseArcWalk}      & false & arc-walk clustering instead of \code{ClusterizeSmooth3D};
                                       recovers track tails \\
\code{SetUseDriftAwareWeights} & false & weight the circle fit by $1/\sigma^2_{xy}(z)$ \\
\code{SetMergeAngleThreshold}  & 45 deg & fragment-merging angular tolerance \\
\bottomrule
\end{tabular}
\end{center}

\note{\textbf{The cluster radius/distance pair is the hidden coupling between PRA and
fitting}: it sets how many measurement points the Kalman filter sees. Too few clusters and
the fit is under-constrained ($\mathrm{ndf}$ collapses); too many and you are fitting
correlated noise. Change it and re-check \code{minClusters} in the fitter.}

\subsection{PSA parameters worth revisiting}

\begin{center}
\begin{tabular}{@{}l l p{9.6cm}@{}}
\toprule
Setter & a1954 & Effect \\
\midrule
\code{SetThreshold}        & 20 & ADC threshold. The single biggest hit-yield knob; lower it to recover faint tracks at the cost of noise (the cleaner then has more work). \\
\code{SetPeakingTime}      & 0.720 & [$\mu$s] GET $\tau$; must match the electronics. \\
\code{SetMaxPeaks}         & 4 & pulses fitted simultaneously per pad --- this is what deconvolves two tracks sharing a pad. \\
\code{SetMinSeparation}    & 4 & [TB] minimum seed spacing; lower resolves closer crossings. \\
\code{SetProminenceOnPrimary} & false & \textbf{leave false.} True rejects broad baseline humps but also kills diffusion-broadened far-drift pulses, collapsing the track in $z$. \\
\code{SetFitChi2Cut}       & 0 (off) & shape gate on the fitted pulse; also enables float-$\tau$. \\
\code{SetMinSignificance}  & 3.0 & keep a pulse only if $A/\sigma_A \ge$ this. \\
\bottomrule
\end{tabular}
\end{center}

\subsection{Cleaner parameters (\code{AtDirDeDxCleaner})}

Constructor defaults: $k=12$, \code{cosSeg}~0.78, \code{cosTan}~0.72, \code{rMax}~32~mm,
\code{qRatio}~0.65, \code{minDeg}~1, \code{smoothQ}~true, \code{qKeep}~150. A hit survives if
it has at least \code{minDeg} neighbours that are both direction- and $dE/dx$-continuous.
Relax \code{cosSeg}/\code{cosTan} or raise \code{rMax} if the cleaner is eating sparse
low-$dE/dx$ tracks; raise \code{minDeg} if noise is surviving.

\subsection{A practical PRA tuning loop}

\begin{enumerate}[leftmargin=*,itemsep=2pt]
\item Pick one representative run and reco only 500--2000 events into a scratch directory,
      varying \emph{one} knob:
\begin{lstlisting}
root -b -q 'pipeline/unpackReco_Be12.C("run_0147",2000,false,"/tmp/scan_mcs15/",...)'
\end{lstlisting}
\item Judge it on the \textbf{PID plane} first (\code{pid/pid_Be12.C}) --- it needs no fit
      and a tighter, more continuous band is the cleanest indicator that clustering
      improved.
\item Then look at track multiplicity per event and at the fraction of events with at least
      one gated track (printed by \code{gate_events_Be12.C}).
\item Only then fit, and judge on the elastic peak: \textbf{FWHM of the $E_x$ g.s.\ peak}
      and the number of surviving tracks. Both are printed by \code{ex_Be12.C} and fitted
      live in the explorers.
\item Keep the scan outputs under distinct \code{outSuffix}/\code{outTag} names so caches
      never overwrite each other.
\end{enumerate}

%===============================================================================
\section{Tuning the fitters}
\label{sec:fit-tuning}
%===============================================================================

\subsection{\code{measSigma} --- and why the $\chi^2$ cut is not portable}

The UKF runs at \code{measSigma = 0.5}~mm and GENFIT at \code{4.0}~mm. Both enter identically
as a per-cluster position sigma, so the $8\times$ difference is $64\times$ in $\chi^2$. This
was measured directly on 73 runs by swapping the two values
(\code{sigma_test_Be12.sh}):

\begin{center}
\begin{tabular}{@{}lrr@{}}
\toprule
& tracks (no $\chi^2$ cut) & kept by $\chi^2/\mathrm{ndf} < 5$ \\
\midrule
UKF at $\sigma = 0.5$    & 32\,120 & 82.8\,\% \\
UKF at $\sigma = 4.0$    & ---     & 99.5\,\% \\
GENFIT at $\sigma = 4.0$ & 31\,578 & 39.7\,\% \\
GENFIT at $\sigma = 0.5$ & ---     & \phantom{0}4.2\,\% \\
\bottomrule
\end{tabular}
\end{center}

\noindent Swapping the sigmas scaled GENFIT's median $\chi^2$ by 63.9 (predicted
$(4/0.5)^2 = 64$) and the UKF's by 40.1 --- less than 64, because the UKF's \code{measSigma}
also enters the filter gain.

\warn{\textbf{A $\chi^2/\mathrm{ndf}$ cut is not portable between fitters or between
\code{measSigma} settings.} With no cut, both fitters return the same tracks to within
1.7\,\%; the shared cut at 5 is the \emph{entire} 26\,604 vs 12\,533 yield gap. Set the cut
per fitter from a percentile of that fitter's own $\chi^2$ distribution, not from a number
copied across.}

\noindent Importantly, \code{measSigma} does \emph{not} move the energy scale: the GENFIT
g.s.\ centroid was $+0.451$~MeV at both $\sigma = 4.0$ and $\sigma = 0.5$ on 31.5\,k tracks
each. Covariance controls \emph{which tracks survive a cut}, not \emph{what energy they get}.

\subsection{UKF knobs}

\begin{longtable}{@{}l l p{8.3cm}@{}}
\toprule
Knob & a1954 & Effect / when to change \\
\midrule
\endhead
\code{measSigma}      & 0.5 mm  & See above. Raise towards the physical cluster resolution if you want a $\chi^2$ that is interpretable; then re-derive the cut. \\
\code{momSigmaFrac}   & 0.1     & Momentum-seed uncertainty. Raise if the seed is poor (short tracks) so the filter is not over-confident. \\
\code{nIter}          & 1       & More iterations help when the seed is far from the solution; costs time linearly. \\
\code{minClusters}    & 10      & Yield/quality trade-off. Lower to keep short backward recoils, at the price of poorly constrained fits. \\
\code{gasDensity}     & 6.5e-5  & CATIMA density. \textbf{Inconsistency to resolve:} the geometry builder computes 6.616e-5 g/cm$^3$ for H$_2$ at 600 torr / 293 K, $\approx 2\,\%$ above this. \\
\code{bFieldSign}     & $-1$    & Experimental handedness. \\
\code{refTrack}       & false   & Do not re-attempt the unscented reference-track port; the UKF cannot match genfit's analytic-Jacobian transport. \\
\code{SetEnableEnergyStraggling} & false & In the macro body. \\
\code{SetEnableMultipleScattering} & --- & In the macro body; worth a test for the lowest-energy protons. \\
\code{SetForceVertexOnBeamAxis} & --- & Vertex-constraint knob in \code{AtFitterUKF} (with \code{SetVertexZBias}), not exposed as a macro argument. \\
\bottomrule
\end{longtable}

\noindent The UKF stores KE at two reference points, which is the cheapest available handle
on energy-loss systematics: \code{GetKinematics()} is the vertex value (back-extrapolated to
the beam axis with the lost energy added back), \code{GetKinematicsXtr()} is the value at the
first cluster. On a1954 the difference is a median of only 0.082~MeV, moving the g.s.\
centroid by $+0.055$~MeV.

\subsection{GENFIT knobs}

\begin{longtable}{@{}l l p{8.3cm}@{}}
\toprule
Knob & Recommended & Effect / when to change \\
\midrule
\endhead
\code{matEffects}   & \code{kTRUE} & \textbf{The important one.} False $\Rightarrow$ constant-momentum helix $\Rightarrow$ KE 30\,\% low below 6~MeV. \\
\code{geoName}      & \code{ATTPC_H600torr} & Material effects are only as good as the density. \\
\code{measSigma}    & 4.0 mm & See Section \ref{sec:fit-tuning}. \\
\code{minIter/maxIter} & 2 / 5 & More iterations for badly seeded tracks. \\
\code{thetaMin}     & 10 ($pp$) / 5 ($pd$) & Rejects near-beam tracks; too high kills forward deuterons. \\
\code{backwardSeedFix} & \code{kFALSE} & Alternative seeding for backward tracks; untested here, worth a scan if backward-angle behaviour still looks odd. \\
\code{pidGate}      & \code{""} & In-fitter Spyral PID gate; redundant when \code{gate_events_Be12.C} already gated. \\
\code{SetELossFile} & (unset) & Loading a proton $dE/dx$ table would close the residual $+16\,\%$ deficit below 2~MeV. \\
\bottomrule
\end{longtable}

\subsection{Drift velocity --- the systematic most likely to be limiting you}
\label{sec:driftv}

The a1954 par file carries \code{DriftVelocity 1.00} cm/$\mu$s for $^{12}$Be. In the sibling
$^{14}$C analysis of the \emph{same experiment}, a scan showed the nominal value (1.30 there)
was clearly too low, and moving to 1.45 improved the resolution by 18\,\% (FWHM
$0.535 \rightarrow 0.441$~MeV) while pulling the g.s.\ from $-0.36$ to $-0.09$~MeV. Three
criteria gave three slightly different optima --- slope$\rightarrow$0 at 1.37, $\mu
\rightarrow 0$ at 1.49, best FWHM at 1.45 --- so quote which one you used.

\textbf{The same scan has not been done for $^{12}$Be and is the highest-value next test.}
The recipe:

\begin{enumerate}[leftmargin=*,itemsep=2pt]
\item Copy \file{parameters/ATTPC.a1954_Be12.par} to
      \file{ATTPC.a1954_Be12_dv110.par} etc., changing only \code{DriftVelocity}.
\item Re-run \emph{reco} for one run per point (drift velocity enters at the PSA
      $z$-calibration, so refitting alone is not enough), into per-point directories.
\item Gate and UKF-fit each point, then build an $E_x$ cache per point with a distinct
      \code{outTag} (remember the leading underscore).
\item Tabulate g.s.\ $\mu$, FWHM and the $dE_x/d\theta_{\rm cm}$ slope against $v_d$, and
      pick your criterion.
\end{enumerate}

\subsection{Other levers on the resolution}

\begin{itemize}[leftmargin=*,itemsep=2pt]
\item \textbf{Beam energy.} $dE_x/dE_{\rm beam} \approx 0.01$~MeV per MeV, so zeroing the
      elastic peak constrains $E_{\rm beam}$ only to $\pm 20$~MeV. Use the \emph{shape} of
      the KE-vs-$\theta_{\rm lab}$ locus (the explorers overlay the kinematic line) or an
      independent measurement.
\item \textbf{Vertex-$z$ dependence.} Resolution is flat in vertex $z$ (a good sign for the
      calibration) but degrades for far vertices; a \code{vz < 600} cut gave FWHM 0.43~MeV
      but cost the weak $2^+$ statistics.
\item \textbf{$E_x$-vs-$\theta$ drift.} If the elastic centroid drifts with angle, that is a
      drift-velocity or KE-scale signature, not a statistical effect. The $^{14}$C
      workspace has \code{pp/hyp_C14.C}, which discriminates the three hypotheses (wrong
      $E_{\rm beam}$ / wrong KE scale / empirical correction) by where it places known
      levels; porting it to $^{12}$Be is straightforward.
\item \textbf{Gates.} Redraw the PID polygon with \code{arclen > 200}; check that the
      deuteron gate is not leaking tritons (this happened in a1975 --- genfit fits them with
      a perfectly good $\chi^2$).
\end{itemize}

%===============================================================================
\section{Known pitfalls}
\label{sec:pitfalls}
%===============================================================================

\begin{enumerate}[leftmargin=*,itemsep=3pt]
\item \textbf{Beam energy is not calibrated.} 155~MeV comes from zeroing the elastic peak,
      which as shown above barely constrains it. Every absolute $E_x$ rides on it. Find the
      logbook value before quoting numbers.
\item \textbf{F: is not auto-mounted} after a reboot: \code{sudo mount -t drvfs F: /mnt/f}.
      Only the raw data and the big reco files live there; all fit outputs are local and
      survive.
\item \textbf{\code{df /} lies inside WSL.} It shows the guest view of a sparse VHDX. When
      the Windows C: drive fills, jobs get killed mid-write while \code{df} still reports
      free space. Watch \code{df /mnt/c} and compact the VHDX to reclaim.
\item \textbf{Kill ROOT jobs at 80\,\% memory} rather than waiting through swap.
\item \textbf{Two runs are timestamp-desynced} (0142, 0212) --- exclude them
      (Section~\ref{sec:tsync}).
\item \textbf{IC amplitude is the positive overshoot}, not $-\min$, and the gate is
      $[500,900]$, not the older $[625,750]$.
\item \textbf{$\theta$ conventions differ} between the gate stage (Spyral polar) and the fit
      output ($\approx 180^\circ$ minus that).
\item \textbf{\code{outTag} needs its leading underscore} --- the cache name is built as
      \code{"proton_kin" + outTag}.
\item \textbf{ROOT Eve / 3D viewers crash under WSL.} Use the browser explorer or static
      PNGs. (Geant4 Qt/OpenGL, by contrast, does work under WSLg.)
\item \textbf{The GENFIT fix propagates.} \code{AtGenfitter} commit \code{5f87e625} affects
      a1975 and PUMA too; anything fitted with \code{matEffects = kTRUE} before it had the
      same silent $\mathrm{ndf} \le 0$ truncation and will change on a re-run.
\end{enumerate}

%===============================================================================
\section{Command cheat-sheet}
%===============================================================================

\begin{lstlisting}
# --- environment ------------------------------------------------------------
export REPO=/home/yassid/fair_install/ATTPCROOTv2-OpenKF
export HERE=$REPO/macro/Unpack_HDF5/a1954_Be12/UKF
source /home/yassid/fair_install/FairSoft/install/bin/thisroot.sh
source $REPO/build/config.sh
export ROOT_INCLUDE_PATH=$REPO/build/include:/home/yassid/fair_install/FairRootInstall/include
sudo mount -t drvfs F: /mnt/f                  # if F: is not mounted
root -b -q $REPO/geometry/ATTPC_H600torr.C     # once per clone

# --- stage 1-2: reco + auxiliary streams ------------------------------------
cd $HERE
./reco_hdb_F.sh      "run_0143 run_0144" 2     # -> /mnt/f/a1954_Be12_reco_hdb/
./frib_F_batch.sh    "run_0143 run_0144" 4     # -> ~/a1954_Be12_reco_hdb_slim/*_FRIB.root
./slim_cache_batch.sh "run_0143 run_0144" 4    # -> ~/a1954_Be12_reco_hdb_slim/*_slim.root

# --- stage 3: checks and gates ----------------------------------------------
root -b -q -l "$HERE/pid/tsync_check.C(\"run_0143\")"
root -b -q -l "$HERE/pid/ic_Be12.C(\"run_0143\")"
root -l      "$HERE/pid/draw_gate_Be12.C(\"proton_12Be\")"     # interactive

# --- stage 4-5: gate + fit ---------------------------------------------------
./fitpipe_Be12.sh    "run_0143 run_0144" 4     # (p,p'): gate -> UKF -> GENFIT
./fitpipe_pd_Be12.sh "run_0143 run_0144" 4 0   # (p,d) : gate -> UKF(deuteron)
./genfit_pd_Be12.sh  "run_0143 run_0144" 4     # (p,d) : GENFIT

# --- stage 6: physics --------------------------------------------------------
root -b -q -l "$HERE/pp/ex_Be12.C(\"run_0143,run_0144\",\"~/a1954_Be12_genfix/\",155,5,\
  \"_pp_genfit\",1.007825,12.026921,\"12Be(p,p')\",\"genfit\")"
$HERE/pp/open_explorer.sh pp
\end{lstlisting}

\subsection*{Output directories}

\begin{center}
\begin{tabular}{@{}ll@{}}
\toprule
\file{/mnt/f/a1954_Be12_reco_hdb/}      & full reco (big) \\
\file{~/a1954_Be12_reco_hdb_slim/}      & slim \code{AtPatternEvent} cache $+$ FRIB \\
\file{~/a1954_Be12_fit/in/}             & gated fit inputs (reusable) \\
\file{~/a1954_Be12_fit/}                & $(p,p')$ UKF $+$ pre-fix GENFIT \\
\file{~/a1954_Be12_genfix/}             & $(p,p')$ fixed GENFIT \\
\file{~/a1954_Be12_fit_pd/}             & $(p,d)$ UKF \\
\file{~/a1954_Be12_genfix_pd/}          & $(p,d)$ fixed GENFIT \\
\file{~/a1954_Be12_sigtest/}            & sigma-swap covariance test \\
\file{$HERE/pp/plots/}                  & $E_x$ caches and PNGs \\
\bottomrule
\end{tabular}
\end{center}

\subsection*{Further reading in the workspace}

\begin{center}
\begin{tabular}{@{}ll@{}}
\toprule
\file{README.md}                    & the terse version of this guide \\
\file{ANALYSIS_REPORT.md}           & the first end-to-end physics pass \\
\file{pp/FITTER_COVARIANCE_TEST.md} & the full UKF-vs-GENFIT investigation \\
\file{pid/PID_COMPARISON.md}        & Spyral vs ATTPCROOT PID \\
\bottomrule
\end{tabular}
\end{center}

\end{document}
